\documentclass[11pt,a4paper]{article}
\usepackage[margin=27mm]{geometry}
\usepackage{fontspec,unicode-math,amsmath,amsthm,booktabs,longtable,hyperref}
\setmainfont{STIXGeneral}\setmathfont{STIX Two Math}\setmonofont{Menlo}
\newtheorem{theorem}{Theorem}\newtheorem{lemma}[theorem]{Lemma}\newtheorem{proposition}[theorem]{Proposition}
\newcommand{\Aut}{\operatorname{Aut}}\newcommand{\Stab}{\operatorname{Stab}}\newcommand{\cl}{\mathcal C}
\title{The complete primitive Mathieu census and its source-table corrections}
\author{Complete corrected application}
\date{}
\begin{document}\maketitle
\begin{abstract}
We recover the 81 rank-two and 431 rank-three geometries in the classification of Leemans and Rowley, using complete maximal-subgroup orbits, an exact chamber-completion test and a separating geometric signature. Every class has an explicit permutation witness and an exact Borel-group identifier. We also reconcile all 431 physical rows of their table without identifying its unavailable historical witness numbering. The later source permits complementary endecad shadows independently; we distinguish that convention from the canonical convention of its predecessor and distinguish two conflicting definitions of the predecessor's pair subscripts. The classification is conditional on the stated complete finite-group catalogue and exact algorithms. No formal verification or historical priority is claimed.
\end{abstract}
This application uses the whole common paper \emph{Coset recognition and exhaustive maximal-subgroup enumeration}. Its definitions, distinguished-group equivalence and Common Lemma~1, Theorem~2 and Proposition~3 are imported together. All group-specific arguments, source qualifications and complete tables remain in this application and its two table supplements. We preserve the predecessor's theorem numbering.
\setcounter{theorem}{3}
\setcounter{section}{1}
\section{The complete Mathieu census}
\begin{theorem}[Leemans--Rowley classification, corrected table edition]\label{thm:census}
For $G=M_{22}$ there are exactly 81 rank-two and 431 rank-three flag-transitive residually connected geometries with maximal proper object stabilizers, under the preceding equivalence. The complete tables accompanying this paper specify every class by a separating signature and its Borel SmallGroups identifier; the permutation witness file constructs every class. The claimed numerical classification is unchanged. The accompanying reconciliation gives a disposition to each of the 431 physical source rows; it is not a recovery of the intended historical \texttt{geo[j]} list.
\end{theorem}
\begin{proof}
Use the complete TomLib $M_{22}$ subgroup catalogue, with 156 conjugacy classes, and its named degree-24 group model. The code compares the advertised maximal indices with the maximal elements of the entire supplied containment relation, constructs all conjugates of the eight representatives, and checks that the resulting 2,300 subgroups are distinct. Catalogue completeness and the named-group identification are retained inputs; this internal comparison alone would not establish them.

The normalizer in the symmetric group has order $1\,774\,080$ and its centralizer of $G$ has order two, from the two fixed points. Its conjugation image therefore has order $887\,040$, compared with the separately calculated order of $\Aut(G)$. Exact action and transporter operations implement Common Proposition~3. They find 81 pair classes, test 11,866 augmented triples, accept 1,193 before final fusion, and return 431 classes. The anticipated answer and the source table are not inputs to this enumeration. Common Lemma~1 proves that every source geometry is represented; Common Theorem~2 proves the converse.

A separate validator checks all witnesses, exact pairwise inequivalence within invariant buckets, all 1,293 chamber-completion equations, all 1,293 residue-generation equations and all 1,293 panel indices. These existence checks supplement the exhaustive traversal; they do not replace it. The complete finite programs and their whole catalogue/software premises are part of the proof. The coordinate and Borel assertions follow from the next two propositions and the finite separating check.
\end{proof}

\begin{proposition}[Geometric shadows]\label{prop:shadows}
In the source's order, the eight maximal classes have the following object model. The final family has two complementary shadows for each subgroup; every other listed shadow determines its stabilizer uniquely.
\[
\begin{array}{c|rrrrrrrr}
\text{type}&1&2&3&4&5&6&7&8\\\hline
|H|&20160&5760&2520&2520&1344&1920&720&660\\
|X|&1&6&7&7&8&2&12&11\\
\text{shadows}&22&77&176&176&330&231&616&1344
\end{array}
\]
They are points, hexads, the two heptad families, octads, duads, dodecads and endecads in the source's $M_{24}$ model.
\end{proposition}
\begin{proof}
Let $\Lambda$ be the 22 moved points and $p,q$ the two fixed points. For every maximal subgroup, the program finds its orbits on $\Lambda$ with the indicated size. There is exactly one except in type 8, where there are two complementary orbits. The set stabilizer of each chosen orbit is checked to equal the original subgroup. Equality of shadows therefore implies equality of stabilizers, and the asserted multiplicities prove the converse statement.

Adjoin both $p,q$ to each hexad, $p$ to one heptad family, $q$ to the other, and retain the octads. The resulting 759 blocks contain each of the $\binom{24}{5}=42\,504$ five-subsets exactly once. Thus they form an actual Steiner system $S(5,8,24)$. We retain its classical uniqueness, its identification with the Curtis MOG and $\Aut(S)=M_{24}$, with pointwise two-point stabilizer of order $443\,520$. Our named $G$ preserves the blocks, fixes both points and has precisely that order, giving the required identification.

Symmetric differences of octads meeting in two points give 2,576 dodecads. The 616 contained in $\Lambda$ are exactly the seventh family; deleting the sole fixed point from a dodecad containing exactly one gives precisely the 1,344 endecads. All equalities are checked as sets, not just as counts. Simplicity of $M_{22}$ makes the action on each proper-subgroup coset set faithful. The abstract maximal-subgroup names in the source's table retain their classical Atlas identification; the computed orbit classification does not infer abstract isomorphism from order alone.
\end{proof}

Fix a distinguished point called $\infty$. For a type-8 subgroup, the \emph{canonical} shadow is the unique $X$ for which $X\cup\{\infty\}$ is a dodecad. For a triple take all orderings with nondecreasing type sequence and both choices of $\infty$, interchanging types 3 and 4 simultaneously. Its signatures are
\begin{equation}\label{eq:signature}
(t_1,t_2,t_3,|X_1\cap X_2|,|X_1\cap X_3|,|X_2\cap X_3|,
 |H_{12}|,|H_{13}|,|H_{23}|).
\end{equation}
For pairs use $(t_1,t_2,|X_1\cap X_2|,|H_{12}|)$. The broader source convention also permits each endecad to be complemented independently. A one-sided complement changes $|X\cap Y|$ to $|Y|-|X\cap Y|$; complementing both changes it to $22-|X|-|Y|+|X\cap Y|$.

\begin{proposition}[Complete separating coordinates]\label{prop:coordinates}
The canonical signature families of the 431 triples are disjoint, as are the broader signature families. The same holds for the 81 pairs. There are respectively 801 and 976 triple signatures and 104 and 115 pair signatures. The least canonical signature gives a complete new table convention.
\end{proposition}
\begin{proof}
Every signature is calculated from the actual stabilizer shadows and subgroup intersections. For each family the finite certificate builds the inverse map from signatures to classes and rejects if any signature belongs to two classes. Every witness has a nonempty family, so disjointness proves separation.

These families are invariants of the specified equivalence. An automorphism of $G$ has a lift in the degree-24 normalizer. On the moved points it sends each shadow to the shadow of the conjugate subgroup. It either preserves both heptad classes or interchanges them; choose its extension on $p,q$ to do the corresponding thing. The definition of the reconstructed octads then shows that this extension preserves all blocks and all dodecads. It preserves the canonical endecad convention or interchanges the two global orientations. The broader family additionally includes the independent complement choices. Type reordering was included from the start. Thus separation is a genuine class invariant, not an artifact of the order of traversal.
\end{proof}

\section{Literal source labels and corrected tables}
The nine-page 2010 primary has 431 physical rows but repeats Number 61 and omits Number 62. We index rows consecutively, reading each printed line left then right; these are distinct from our witness identifiers. The complete row table preserves its page, label, Number and Borel text. Neither numbering is silently substituted for the other.

The 2003 predecessor explicitly chooses canonical endecads on printed page 109. The 2010 source defines endecads by either distinguished-point extension and expressly allows complements on pages 2--3. Accordingly our primary row diagnostic permits independent complements. Restricting to the older convention is a separate sensitivity: rows 178, 280 and 394 acquire compatible witnesses under the later convention. They are not errors proved by canonical-parity arguments.

Pair subscripts require a further distinction. The predecessor's printed page 112 distinguishes eleven pairs of subscript values by pair-intersection orders. Its page 114 gives 22 explicit pairs of objects. These are not entirely consistent. We give separate candidate columns for the order-table and literal-object readings, recording undefined object labels rather than guessing replacements. For ordinary named Borel groups the test uses exact SmallGroups identifiers. A juxtaposition $p^aH$ is tested only for a normal elementary abelian subgroup of order $p^a$ and the named quotient; it does not assert an unprinted action or splitting. The replacement table itself gives exact group identifiers throughout.

\begin{lemma}[All compatible labelled realizations]\label{lem:codes}
The twenty unambiguous octads transcribed from the predecessor span a self-orthogonal binary code $C_0$ of dimension ten. Exactly four extended Golay codes in those labelled coordinates contain $C_0$; all give the same outcomes for the 22 literal pairs and three colon configurations used here.
\end{lemma}
\begin{proof}
Binary elimination checks dimension ten and every scalar product of basis vectors. The orthogonal space has dimension fourteen. A self-dual dimension-twelve extension lies in $C_0^\perp$ and corresponds to a two-dimensional subspace of the four-dimensional quotient. Enumerating all pairs of distinct nonzero quotient vectors, then deduplicating their spans, gives all 35 such planes. The all-one vector belongs to $C_0$. Retain exactly the planes whose full preimage has word weights in $\{0,8,12,16,24\}$; minimum distance eight and double evenness force precisely these weights, since weight twenty would give a weight-four complement. There are four retained planes, each with 4,096 words and distribution $1,759,2576,759,1$. Every five-subset lies in precisely one weight-eight support. Double evenness implies self-orthogonality by the weight identity, and dimension twelve gives self-duality.

The classical identification of the Curtis MOG with the extended Golay code is retained. Every realization satisfying the literal seed octads therefore occurs in this enumeration. For each completion, a coloured incidence-graph isomorphism to the constructed Steiner system fixes the two specified points; all 759 block images are then verified directly. Actual set stabilizers, local generation and chamber-completion counts are computed in every completion. The equality of their reported outcomes proves independence from any guessed coordinate labelling.
\end{proof}

The literal pair called $M_{56}(0_2)$ has intersection-stabilizer order 192, whereas the order table prints 16. The second eleven-set printed for $M_{88}(5_1)$ is not the required endecad and has stabilizer order 12 rather than 660. The label $e_4$ used in another pair is not assigned in the named-object list: its fourth displayed endecad is labelled $e_3$ again. We record the fourth-display interpretation provisionally, with its status visible on every affected row.

For the three colon labels, the literal configurations give the following complete test. All local residue-generation orders equal the required maximal orders and each Borel has order two.
\[
\begin{array}{c|c|c|c}
\text{label}&(|H_{12}|,|H_{13}|,|H_{23}|)&
 \text{three completion counts / required}&\text{conclusion}\\\hline
M_{336}(3,0,1:2)&(24,24,24)&(44/12,44/12,44/12)&\text{not flag-transitive}\\
M_{366}(0,0,0_1:2)&(24,24,12)&(33/6,66/12,66/12)&\text{not flag-transitive}\\
M_{378}(2,2,4:2)&(20,12,6)&(3/3,6/6,10/10)&\text{class 338}
\end{array}
\]
The final class is identified by an exact automorphism transporter, not just its signature. These tests concern the displayed configurations, not conjectured intended replacements. Independently, $M_{133}(0,3,1)$ asks a point to meet a heptad in three points and is impossible. The two literally identical labels in rows 370 and 375, and the identical labels with different Borel descriptions in rows 105 and 106, are preserved explicitly.

The five rank-two duplicate labels identified in 2010 correspond to complementary shadows: the intersection values change as
\[
 M_{18}:1\leftrightarrow0,\quad M_{28}:5\leftrightarrow1,\quad
 M_{58}:6\leftrightarrow2,\quad M_{68}:2\leftrightarrow0,\quad
 M_{78}:8\leftrightarrow4.
\]
Complementation fixes the endecad stabilizer, so these pairs define the same coset geometry. The independent 81-class table supplies every surviving rank-two class. This does not silently certify the conflicting explicit definitions on the predecessor's page 114.

The complete corrected result is the exhaustive class table together with the literal-row dispositions. An incompatible or undefined source label is not used as a premise and is not assigned a guessed historical number. Recovering an author's intended unavailable witness list is outside this claim. Thus the 431-class theorem is established, while the stronger assertion that all printed labels are simultaneously correct is explicitly replaced by a complete corrected table.

\section{Attribution and exact proof boundary}
Leemans--Rowley proved the stated numerical classification and supplied the geometric notation and source table. Kılıç--Rowley supplied the predecessor's MOG coordinates and refined labels. The coset realization, connectedness-by-generation argument and stabilizer-orbit augmentation are classical. The previous project phases 164 and 173 already established and independently reviewed the exhaustive census and geometric shadows; phase 180 supplied the labelled Golay-completion and canonical-signature work, then only an author candidate. This edition reuses that work, supplies the full written reductions here, extends the signature table to the actual 2010 complement convention and separates the conflicting pair interpretations. It does not rely on the old umbrella manuscript, on the special $J_1\times\operatorname{PSL}_2(19)$ theorem, or on a missing downloaded geometry list.

The required finite premises are the complete GAP/TomLib subgroup and named-group data, exact permutation-group and graph algorithms, and SmallGroups identifications. The classical Steiner, Golay, Mathieu and Atlas identifications remain explicit outside inputs. All mathematical programs, source-coordinate data, permutation witnesses and generated complete tables are bound and supplied; the replay driver is separate. There is no Lean proof, no newly solved conjecture, and no compression factor in this packet. A separate independent review is required for admission of this corrected complete source contract.

\begin{thebibliography}{9}
\bibitem{LR} D. Leemans and P. Rowley, \emph{Rank three residually connected geometries for $M_{22}$, revisited}, Electronic Journal of Combinatorics 17 (2010), N4, 9 pp.
\bibitem{KR} N. Kılıç and P. Rowley, \emph{On rank 2 and rank 3 residually connected geometries for $M_{22}$}, Note di Matematica 22 (2003), 107--154. The specific defining pages are 108--114 and 143, 149--150; no complete new recovery of this 48-page predecessor is claimed.
\bibitem{Curtis} R. T. Curtis, \emph{A new combinatorial approach to $M_{24}$}, Math. Proc. Cambridge Philos. Soc. 79 (1976), 25--42. Classical identification retained through the source's explicit model.
\bibitem{Atlas} J. H. Conway et al., \emph{Atlas of Finite Groups}, Oxford University Press, 1985. Named Mathieu and maximal-subgroup identifications retained.
\end{thebibliography}
\end{document}
